\section{Життєвий цикл криптографічного ключа}

Спрощена схема життєвого циклу криптографічного ключа охоплює такі послідовні етапи: генерація
$\rightarrow$ зберігання $\rightarrow$ використання $\rightarrow$ ротація/заміна $\rightarrow$ відкликання
або знищення. Для забезпечення безпеки кожен із цих етапів вимагає врахування специфічних загроз та
дотримання суворих вимог.

На етапі генерації ключа основною загрозою є використання передбачуваних алгоритмів або недостатньої
ентропії (наприклад, застосування звичайних некриптографічних PRNG чи слабкого початкового значення
\textit{seed}), що дозволяє зловмиснику відтворити ключовий матеріал шляхом повного перебору. Вимогою
безпеки тут є обов'язкове використання апаратно підкріплених криптографічно стійких генераторів
випадкових чисел (CSPRNG) у захищеному середовищі~\cite{SP800-90A}.

На етапі зберігання ключа виникають загрози несанкціонованого доступу до файлової системи,
конфігураційних файлів чи вихідного коду програми, де ключі можуть зберігатися у відкритому вигляді, а
також ризики витоку через дампи пам'яті. Вимогою безпеки є розміщення ключів виключно в ізольованих
шифрованих сховищах з використанням обмеження прав доступу та заборони експорту
(\textit{non-exportable keys})~\cite{SP800-57}.

Під час використання ключа основною загрозою є перехоплення секретних даних в оперативній пам'яті,
атаки за побічними каналами (часові атаки або аналіз кешу процесора), а також випадковий запис ключа до
системних журналів (логів) під час діагностики. Вимогами безпеки на цьому етапі є виконання
криптографічних операцій із застосуванням алгоритмів із постійним часом виконання
(\textit{constant-time implementation}), ретельна фільтрація журналів та оперативне очищення ділянок пам'яті
після завершення обчислень~\cite{SP800-57}.

На етапі ротації або заміни ключа ключовою загрозою є відсутність регулярного оновлення ключового
матеріалу, що призводить до надмірного накопичення зашифрованих даних під захистом одного ключа та
збільшує масштаб можливої компрометації. Вимогою безпеки є впровадження автоматизованої політики
планової ротації за розкладом або після обробки визначеного обсягу інформації з підтримкою
контрольованого перехідного періоду~\cite{SP800-57}.

На завершальному етапі відкликання або знищення основними загрозами є неповне або поверхневе
видалення ключа, внаслідок якого залишки ключової інформації можуть бути відновлені зі стаціонарних
носіїв чи кешу, а також затримка блокування компрометованого ключа системою. Вимогами безпеки тут є
реалізація гарантованого безповоротного затирання пам'яті (процедури \textit{zeroization} або
\textit{explicit\_bzero}) та миттєве внесення скомпрометованих ключів до чорних списків відкликання~\cite{FIPS140-3}.
